% RouteSAM 中文理解稿
% 建议使用 XeLaTeX 编译；本稿与英文投稿稿共享公式、图片和参考文献。
\documentclass[UTF8,a4paper,10pt]{ctexart}

\usepackage[left=2.1cm,right=2.1cm,top=2.2cm,bottom=2.2cm]{geometry}
\usepackage{amsmath,amssymb,bm}
\usepackage{graphicx}
\usepackage{float}
\usepackage[authoryear,round]{natbib}
\usepackage[dvipsnames]{xcolor}
\usepackage[colorlinks=true,citecolor=TealBlue,linkcolor=blue,urlcolor=blue]{hyperref}
\usepackage{booktabs}
\usepackage{microtype}

\setlength{\parindent}{2em}
\setlength{\parskip}{0.25em}
\linespread{1.15}

\title{RouteSAM：通过跨图像语义路径实现半监督二维医学图像分割}
\author{作者信息待补充}
\date{}

\begin{document}
\maketitle

\begin{abstract}
医学图像的高昂标注成本使半监督学习成为医学图像分割中的重要研究方向。然而，当有限的标注数据难以充分覆盖多样化的病灶表征时，半监督分割模型往往会出现明显的性能下降。尤其是在标注数据与未标注数据之间存在较大语义差异的情况下，有限的标注信息难以稳定迁移至未标注图像，从而影响所生成伪监督的可靠性。
为此，我们提出 RouteSAM，一种基于语义路由的渐进式监督传播方法。为使有限监督覆盖尽可能广泛的特征空间，RouteSAM 首先从训练数据中筛选具有较高特征覆盖性的样本，作为后续监督传播的起点。进一步地，针对未标注图像，RouteSAM 在其与已标注图像之间构建由多个特征相关图像组成的传播路径，使监督信息沿语义邻近的图像逐步传递，从而避免在特征差异较大的图像之间进行困难的直接迁移。
在此基础上，我们进一步设计了一种多阶段的通用模型--专用模型协同机制，利用 SAM3 的通用分割能力与任务特定 U-Net 的医学域知识，通过伪监督的迭代筛选与模型间的协同适配，逐步提高监督信息的可靠性并增强模型对目标医学任务的适应能力。
在 Kvasir-SEG、BUSI 和 ISIC2018 数据集上的实验结果表明，RouteSAM 在极少量标注数据下能够有效提升半监督医学图像分割性能，\textbf{[最终定量结果待补充]}。

\noindent\textbf{关键词：}半监督医学图像分割；Segment Anything Model；跨图像传播；语义路径；伪标签学习
\end{abstract}

\section{Introduction}

医学图像分割通过对病灶、器官及其他关键解剖结构进行精确分割，为计算机辅助诊断、疾病定量分析和治疗规划等任务提供重要的空间信息。近年来，深度学习方法显著提升了医学图像分割性能，但其有效训练通常依赖大量高质量的像素级标注。不同于自然图像，医学图像的精细标注往往需要具备专业知识的医师参与，标注过程耗时且成本较高，使得大规模、高质量标注数据难以获得~\cite{hesamian2019deep,tajbakhsh2020embracing}。为降低对人工标注的依赖，半监督医学图像分割（Semi-Supervised Medical Image Segmentation, SSMIS）通过联合利用少量有标注数据和大量无标注数据，已成为降低医学图像标注需求的重要研究方向~\cite{jiao2024learning}。

现有半监督医学图像分割方法主要通过伪标签学习和一致性约束从无标注数据中获取额外监督信息。一类方法利用模型对无标注样本的预测构造伪标签，并通过教师--学生或多网络协同机制不断更新监督信号；另一类方法则通过对输入、特征或模型施加不同扰动，约束同一样本在不同条件下保持预测一致性。近年来，Cross Teaching、MC-Net+ 和 BCP 等方法分别从异构网络互监督、不确定区域一致性学习以及有标注与无标注样本间的数据混合等角度提高了无标注数据的利用效率~\cite{luo2022crossteaching,wu2022mcnet,bai2023bcp}。然而，当标注样本进一步减少时，模型能够获得的可靠任务知识也随之受限，由模型自身预测产生的伪标签更容易受到预测误差和不确定性的影响。因此，在极低标注条件下，如何为大量无标注图像获得稳定、可靠的伪监督，仍然是半监督医学图像分割中的关键问题。

可提示分割基础模型的出现为极低标注条件下获取更可靠的监督信息提供了新的可能。以 Segment Anything Model（SAM）为代表的基础分割模型通过大规模数据预训练获得了较强的通用视觉表征与提示驱动分割能力，并能够借助点、框或掩码等提示完成目标区域分割~\cite{kirillov2023segment}。尽管自然图像与医学图像之间存在明显的领域差异（domain gap），使得 SAM 在不同医学影像任务上的直接应用并不稳定~\cite{mazurowski2023samexperiment,ma2024medsam}，其大规模预训练所获得的通用分割先验仍具有较高的利用价值。因此，近年来一些研究开始将 SAM 引入半监督医学图像分割，通过任务特定模型生成提示，并利用 SAM 的预测构造或优化无标注样本的伪监督。例如，SemiSAM 利用任务特定模型提供定位信息并生成 SAM 的输入提示，SemiSAM+ 进一步构建了通用模型与任务特定模型之间的协同学习机制，而 SAMatch 则借助 SAM 对已有预测进行进一步优化~\cite{zhang2023semisam,zhang2025semisamplus,xu2024samatch}。这些工作表明，基础模型所包含的通用分割先验能够为有限标注条件下的无标注数据学习提供额外的监督来源。

尽管上述方法能够借助基础模型提高伪监督质量，其关注点仍主要集中在如何获得更准确的单图像预测或伪标签。与此同时，已有研究表明，医学图像之间存在可利用的跨样本语义关系。例如，ACTION 通过建模全局语义关系和局部解剖信息提升半监督表示学习，说明训练集中不同样本之间并非彼此独立~\cite{you2022action}。另一方面，YoloSeg 进一步利用 SAM2 将单个标注样本的分割信息传播至无标注数据，在极低标注条件下证明了跨图像传播（cross-image propagation）的应用潜力~\cite{zhang2026yoloseg}。然而，这类方法主要采用从标注样本到目标图像的直接传播（direct propagation）。当两幅图像在病灶形态、尺度或外观上存在较大的语义差异（semantic discrepancy）时，一次直接传播仍需跨越较大的样本差异，因而可能产生不稳定的伪监督。由此产生一个值得进一步研究的问题：能否利用无标注数据内部的跨样本关系，在标注样本与困难目标之间引入语义相关的中间图像，从而将一次困难的直接传播分解为多个更局部的传播过程？

基于这一问题，我们进一步考虑如何利用无标注数据内部的样本关系来组织有限监督的传播。如图~\ref{fig:motivation} 所示，我们将无标注训练集视为一个潜在的语义过渡空间（semantic transition space），并提出 \textbf{RouteSAM}。在固定标注预算下，RouteSAM 首先从训练数据中选择具有较强分布覆盖能力的标注样本作为 \textbf{Anchor}。对于每个需要生成伪标签的无标注 \textbf{Target}，RouteSAM 从无标注数据中寻找若干语义相关的 \textbf{Bridge}，并依据样本之间的语义关系将其组织为一条有序的语义路径（semantic route）：
\begin{equation}
A \rightarrow B_1 \rightarrow B_2 \rightarrow \cdots \rightarrow B_K \rightarrow T.
\end{equation}
随后，SAM3 沿该语义路径执行渐进式跨图像传播（progressive cross-image propagation），逐步将 Anchor 的分割信息传递至 Target。与直接传播相比，该过程通过 Bridge 将一次跨度较大的传播分解为多个更局部的传播步骤，使无标注图像不仅是等待生成伪标签的 Target，也能够作为其他 Target 语义路径中的 Bridge，参与有限监督的逐步扩展。

由于不同 Anchor 和不同语义路径可能为同一个 Target 产生不同的掩码候选（mask candidates），直接使用所有传播结果仍可能引入不可靠的伪监督。为此，RouteSAM 进一步评估不同掩码候选的质量，并从中选择可靠的路由伪标签（route-derived pseudo-labels）用于后续学习。在此基础上，我们利用少量 Anchor 的真实标注和路由伪标签共同训练任务特定分割模型，并将该模型称为 \textbf{specialist}；SAM3 则作为具有通用分割先验的 \textbf{generalist}。Specialist 从目标医学数据中学习任务特定知识，并进一步为 generalist 的传播结果提供反馈，使通用分割能力与任务特定知识能够形成互补。语义路由和跨图像传播仅用于训练阶段，最终推理由 specialist 独立完成，无需检索 Anchor、构建语义路径或提供人工提示。

为验证 RouteSAM 在极低标注条件下的有效性，我们在 Kvasir-SEG、BUSI 和 ISIC2018 三个具有不同成像特征与病灶类型的二维医学图像分割数据集上开展实验，并在约 $1\%$ 的标注预算下进行统一评估。我们将 RouteSAM 与代表性的半监督医学图像分割方法、SAM 辅助的低标注分割方法以及直接跨图像传播方法进行比较，以考察跨图像监督路由（cross-image supervision routing）在有限监督扩展中的作用。同时，通过对 Anchor 选择、Bridge 引入以及语义路径构建等关键因素进行消融与分析，进一步研究不同路由设计对伪监督质量和最终分割性能的影响。

本文的主要贡献概括如下：
\begin{itemize}

\item
针对极低标注条件下有限监督难以稳定到达语义差异较大的无标注图像的问题，我们提出 \textbf{RouteSAM}。RouteSAM 从跨图像监督路由的角度利用训练集内部的样本关系，使无标注数据能够参与有限监督的逐步扩展。

\item
我们提出 \textbf{Coverage-Aware Anchor Selection}。在固定标注预算下选择具有较强训练分布覆盖能力的样本作为 Anchor，为后续跨图像监督传播提供更具代表性的起点。

\item
我们提出 \textbf{Anchor--Bridge--Target} 语义路由机制。通过检索并组织语义相关的 Bridge 构建有序语义路径，并利用 SAM3 执行渐进式跨图像传播，将困难的直接传播分解为多个更局部的传播步骤。

\item
我们提出面向路由伪监督的可靠学习与 \textbf{generalist--specialist} 协同策略。通过筛选不同语义路径产生的掩码候选获得可靠路由伪标签，并利用 specialist 的任务特定知识进一步改善 generalist 的传播结果。

\end{itemize}
\section{相关工作}

\subsection{半监督医学图像分割}

半监督医学图像分割利用少量像素级标注和大量未标注图像训练模型。一类代表性方法是一致性正则化：同一图像在不同扰动、不同网络或不同视图下应得到一致预测。Mean Teacher 使用滑动平均教师产生稳定目标~\cite{tarvainen2017meanteacher}；Cross Pseudo Supervision 让两个独立初始化的网络交换硬伪标签~\cite{chen2021cps}；Cross Teaching 利用 CNN 与 Transformer 的不同归纳偏置进行互教~\cite{luo2022crossteaching}；MC-Net+ 则通过多个解码器估计不确定区域并施加软一致性~\cite{wu2022mcnet}。

近期工作进一步提高了伪监督的可靠性。BCP 通过双向复制粘贴减小标注数据与未标注数据之间的经验分布差异~\cite{bai2023bcp}；ABD 使用置信度引导的双向区域置换控制多种扰动~\cite{chi2024abd}；DyCON 对不确定体素动态加权，并结合一致性和对比学习~\cite{assefa2025dycon}。这些方法主要通过预测一致性、扰动或局部样本混合来利用未标注数据。RouteSAM 与其区别在于，它不只改进单张图像的伪标签，而是显式构造一条由多张图像组成的监督传播路径。

\subsection{基础分割模型辅助的少标注学习}

SAM 的提示式分割能力使基础模型能够为少标注医学任务提供额外先验。SemiSAM 将冻结的 SAM 作为传统一致性框架的辅助分支：任务模型产生提示，SAM 输出额外伪监督~\cite{zhang2023semisam}。SemiSAM+ 将这种设计推广到一个或多个冻结的通用模型，并加入置信度感知的通用模型--任务模型协作~\cite{zhang2025semisamplus}。SAMatch 结合弱到强的一致性学习和任务适配后的 SAM，由教师自动生成提示，再由 SAM 返回精化掩码~\cite{xu2024samatch}。

另一些方法直接利用未标注数据适配 SAM。CPC-SAM 使用两个 SAM 解码器相互产生提示与监督，并约束不同提示位置下的预测一致性~\cite{miao2024cpcsam}。SC-SAM 让 U-Net 为 SAM 提供点提示和伪标签，同时由 SAM 反向约束 U-Net~\cite{vu2026scsam}。CPPS-SAM 也通过交叉提示与交叉伪监督，在 CNN 专家和 SAM 通用模型之间双向传递知识~\cite{zhang2026cppssam}。这些方法证明了通用模型和任务模型协作的价值，但协作主要发生在同一图像内部。RouteSAM 关注的是另一个维度：监督应当如何在一组相关图像之间被组织和传播。

\subsection{跨图像关系建模与掩码传播}

一些半监督方法已经利用了数据集中的跨样本关系。ACTION 学习全局语义关系和局部解剖对比~\cite{you2022action}；DACL 使用密度感知的邻居图将稀疏特征拉向高密度类别区域~\cite{tang2024dacl}；GraphCL 将实例图和图聚类引入教师--学生框架~\cite{wang2024graphcl}。这些研究说明未标注数据包含有价值的数据集级结构，但跨样本关系主要用于改善特征空间，并没有直接回答掩码应当沿哪些图像传播。

参考图像驱动的方法与本文更为接近。PerSAM 和 Matcher 根据一张参考图像及其掩码个性化或提示 SAM~\cite{zhang2023persam,liu2023matcher}；OP-SAM 通过支持图像与查询图像之间的相关性和迭代提示，将单个标注样本的先验迁移到息肉图像~\cite{mao2025opsam}；YoloSeg 使用 SAM2 将单个标注传播到未标注图像，并通过多视图共识训练任务模型~\cite{zhang2026yoloseg}。RouteSAM 的关键区别在于：它不仅决定从哪张标注图像出发，还在起点与目标之间主动插入并排序未标注桥接图像，把未标注池从被动的预测对象转化为可用于渐进监督迁移的过渡空间。

\section{方法}
\label{sec:method-zh}

\subsection{问题定义与整体思路}

设训练图像集合为 $\mathcal{D}_{\mathrm{tr}}=\{x_i\}_{i=1}^{N}$，固定的验证集和测试集分别为 $\mathcal{D}_{\mathrm{val}}$ 与 $\mathcal{D}_{\mathrm{te}}$。在约 1\% 的标注预算下，我们从训练集中选择
\begin{equation}
K=\max\left(1,\operatorname{round}(0.01N)\right)
\end{equation}
张图像作为锚点并获取人工掩码，其余训练图像构成未标注池 $\mathcal{U}$。验证集标签只用于拟合和选择质量估计模块；所有测试预测冻结之前，测试标签始终不可见。

RouteSAM 分为三个阶段。阶段 0 在不知道任何训练掩码的前提下，从训练图像中选择尽可能覆盖数据分布的锚点。阶段 1 针对每个目标检索合适的锚点和桥接图像，构造不同深度的传播路径，并使用冻结的 SAM3 生成与路径相关的目标候选掩码。阶段 2 用这些伪标签训练单图像任务模型，再由任务模型审核传播结果并辅助 SAM3 适配；适配后的 SAM3 沿固定路径重新传播，最终训练部署模型。

\subsection{覆盖感知的锚点选择}
\label{sec:anchor-selection-zh}

锚点不仅是一张标注样本，还是许多传播路径的起点。如果少量锚点都集中在相似外观上，那么其他类型的目标只能从很不匹配的起点获得监督。因此，我们不进行独立随机抽样，而是选择能够共同覆盖训练分布的图像。整个选择过程只读取 RGB 图像；锚点列表完全冻结后才打开对应掩码。

\paragraph{全局描述。}
每张训练图像以 $1008\times1008$ 的分辨率输入冻结的 SAM3 图像骨干。设其 patch token 为 $\mathbf{P}_i=\{\mathbf{p}_{i,r}\}_{r=1}^{72\times72}$。先归一化每个 token，再进行平均，得到全局描述：
\begin{equation}
\mathbf{g}_i=\operatorname{norm}\left(\frac{1}{|\mathbf{P}_i|}\sum_r
\operatorname{norm}(\mathbf{p}_{i,r})\right),
\end{equation}
其中 $\operatorname{norm}(\mathbf{v})=\mathbf{v}/\lVert\mathbf{v}\rVert_2$。

\paragraph{局部描述。}
仅使用全局平均可能忽略小目标和局部纹理，因此我们从每张图像均匀采样 $8\times8$ 个 patch token，并在训练集上拟合包含 64 个聚类中心的 MiniBatch K-Means 视觉词典。记 $h_{i,c}$ 为图像 $i$ 中视觉词 $c$ 的归一化频数，并定义
\begin{equation}
\operatorname{idf}_c=\log\frac{N+1}{\operatorname{df}_c+1}+1.
\end{equation}
局部描述为经过 Hellinger 变换并归一化的直方图：
\begin{equation}
\mathbf{l}_i=\operatorname{norm}\left(\sqrt{\mathbf{h}_i\odot\mathbf{idf}}\right).
\end{equation}
综合全局与局部外观，图像相似度定义为
\begin{equation}
S(i,j)=\frac{1}{2}\left[\mathbf{g}_i^{\top}\mathbf{g}_j\right]_+
+\frac{1}{2}\left[\mathbf{l}_i^{\top}\mathbf{l}_j\right]_+,
\end{equation}
其中 $[z]_+=\max(z,0)$。

\paragraph{贪心覆盖。}
对于锚点集合 $\mathcal{A}$，使用设施选址形式的覆盖目标：
\begin{equation}
F(\mathcal{A})=\frac{1}{N}\sum_{i=1}^{N}\max_{a\in\mathcal{A}}S(i,a).
\end{equation}
从空集合开始，每次加入使 $F$ 的边际增益最大的图像，直到选出 $K$ 个锚点。完全重复的图像文件不作为重复候选，但仍保留在覆盖目标中。相同得分按照固定样本编号打破平局。最终锚点列表及其哈希在读取掩码之前冻结。

\subsection{面向目标的锚点检索}
\label{sec:target-pooling-zh}

覆盖选择回答“哪些图像值得被标注”，但不同目标需要的起点并不相同。阶段 1 因而使用目标感知的检索方法，从锚点集合中找到与当前目标最相关的一个或两个锚点。

检索与传播使用独立的 $256\times256$ 表示。SAM3 骨干为每张图像产生 $18\times18$ 个 1024 维 patch token，记为 $\{\mathbf{x}_{i,j}\}_{j=1}^{324}$。对于锚点 $a$，把其掩码缩放到 patch 网格，并平均前景位置 $\Omega_a$ 内的 token，得到前景原型
\begin{equation}
\mathbf{p}_a=\operatorname{norm}\left(
\frac{1}{|\Omega_a|}\sum_{j\in\Omega_a}\operatorname{norm}(\mathbf{x}_{a,j})
\right).
\end{equation}
若缩放后前景为空，则退化为整幅图像的 patch 均值。

对于目标 $t$，将锚点前景原型与目标各 patch 的相似度转化为注意力权重：
\begin{equation}
w_{a,t,j}=\frac{\exp(\beta\mathbf{p}_a^{\top}\mathbf{x}_{t,j})}
{\sum_k\exp(\beta\mathbf{p}_a^{\top}\mathbf{x}_{t,k})},
\qquad \beta=10.
\end{equation}
由此得到锚点条件下的目标描述和 Target Pooling 分数：
\begin{equation}
\mathbf{z}(a,t)=\operatorname{norm}\left(\sum_jw_{a,t,j}\mathbf{x}_{t,j}\right),
\qquad
r(a,t)=\mathbf{p}_a^{\top}\mathbf{z}(a,t).
\end{equation}
默认使用 $r(a,t)$ 排序锚点。我们同时研究中心化分数 $r(a,t)-\mu_a$，其中 $\mu_a$ 是锚点 $a$ 在训练 RGB 图像上的平均 Target Pooling 分数。中心化只作为锚点排序的校准消融，不改变后续路径目标。最终保留排名最高的一个或两个锚点。

\subsection{有序语义路径构造}
\label{sec:route-construction-zh}

对于每个保留的锚点--目标组合，RouteSAM 分别构造包含 0 至 6 张桥接图像的路径：
\begin{equation}
\mathcal{R}_{b}: A\rightarrow B_1\rightarrow\cdots\rightarrow B_b\rightarrow T,
\qquad b\in\{0,\ldots,6\}.
\end{equation}
所有桥接图像只能来自训练划分，目标图像不能出现在自己的桥接链中。相邻图像之间的亲和度由归一化全图 patch 均值描述的余弦相似度计算。对于一条候选路径，我们同时考虑中间节点及目标的 Target Pooling 分数，以及所有相邻边的亲和度，并用下式评分：
\begin{equation}
Q(\mathcal{R})=\left(\min_{v\in\mathcal{V}(\mathcal{R})}v,
\frac{1}{|\mathcal{V}(\mathcal{R})|}\sum_{v\in\mathcal{V}(\mathcal{R})}v\right).
\end{equation}
该二元组按字典序比较：首先最大化路径中最弱的节点或转移，再最大化平均相关性。直观上，它避免某一段严重不匹配的路径，而不是只追求较高的平均分。每种深度分别使用宽度为 32 的束搜索，避免短路径与长路径相互挤占搜索空间。

\subsection{SAM3 渐进传播}
\label{sec:sam3-propagation-zh}

每条路径被视作一段短伪视频，各帧按照式（10）的顺序排列。SAM3 在锚点帧上由人工掩码的紧致外接框初始化，不使用文本提示。模型随后在 $256\times256$ 画布上将目标掩码依次传播经过各个桥接帧，最终目标帧上的二值掩码被保存为路径候选 $\hat{y}_{t,\mathcal{R}}$。

由于目标真值不可见，我们还执行一次反向传播来获得无标签质量信号。具体而言，用预测得到的目标掩码初始化逆序路径，并把返回锚点的预测与已知锚点掩码比较，得到循环一致性 $q_{\mathrm{cycle}}$。该比较只发生在有标注的锚点，因此既能用于未标注训练目标，也能用于独立测试目标。

\subsection{路径候选质量估计}
\label{sec:router-zh}

若保留 $M$ 个锚点，则每个目标共有 $7M$ 个候选掩码。RouteSAM 使用 28 个不依赖目标真值的特征描述候选。7 个基础特征包括路径深度、路径最小与平均亲和度、最终 SAM 分数、候选数量及深度--亲和度交互；另外 21 个传播特征刻画反向传播的一致性与成功状态、掩码面积变化、空掩码与连通分量、质心和外接框漂移、相邻帧掩码一致性、SAM 分数轨迹以及候选数量轨迹。

这些特征标准化后输入 $\alpha=1$ 的 Ridge 回归器，在验证集上预测候选 Dice。五折拟合时，同一验证图像的全部候选始终属于同一组，从而避免候选级信息泄漏。推理时选择预测质量最高的候选，并用固定的锚点排名和路径深度规则打破平局。所有测试选择在读取测试真值之前冻结。同一候选池上的逐目标 Oracle 仅作为诊断上界，不参与部署。

\subsection{任务模型审核与 SAM3 适配}
\label{sec:collaboration-zh}

阶段 1 选出的掩码形成第一轮伪标注集合
\begin{equation}
\widehat{\mathcal{D}}_{U}^{(1)}=
\{(x_t,\hat{y}_t^{(1)},q_t^{(1)})\mid x_t\in\mathcal{U}\},
\end{equation}
其中 $q_t^{(1)}$ 是预测的路径质量。任务专用 U-Net $f_{\theta}$ 使用锚点真值和通过筛选的伪标签训练：
\begin{equation}
\begin{aligned}
\mathcal{L}_{\mathrm{sup}}&=\frac{1}{K}\sum_{a\in\mathcal{A}}
\ell(f_{\theta}(x_a),y_a),\\
\mathcal{L}_{\mathrm{pseudo}}&=\frac{1}{|\widehat{\mathcal{D}}_{U}^{(1)}|}
\sum_t q_t^{(1)}\ell(f_{\theta}(x_t),\hat{y}_t^{(1)}),\\
\mathcal{L}_{\mathrm{student}}&=\mathcal{L}_{\mathrm{sup}}
+\lambda_u\mathcal{L}_{\mathrm{pseudo}},
\end{aligned}
\end{equation}
其中 $\ell$ 由二元交叉熵和 soft Dice 损失组成。

U-Net 随后作为独立的领域审核者。对每个路径候选，我们综合传播质量估计和候选掩码与 U-Net 预测之间的一致性。接受阈值及组合权重只根据验证集选择。其目的不是简单地让两个模型互相确认，而是发现那些沿路径看似自洽、却不符合目标领域外观的传播结果。

通过审核的伪标签和锚点掩码用于参数高效的 SAM3 适配。为了把“模型适配收益”与“重新搜索路径的收益”区分开，适配之前冻结锚点身份、桥接成员、路径顺序和候选选择策略。

% TODO（阶段 2）：自动锚点闭环实现后，冻结审核公式、阈值、PEFT 模块和优化日程。

\subsection{保持路径不变的第二轮传播与最终推理}
\label{sec:repropagation-zh}

将适配后的 SAM3 应用于相同的冻结路径，得到第二轮候选 $\hat{y}_{t,\mathcal{R}}^{(2)}$。随后按照预先在验证集确定的策略再次执行质量估计和任务模型审核，构造 $\widehat{\mathcal{D}}_{U}^{(2)}$。最终 U-Net 使用锚点真值和第二轮伪监督训练，目标函数仍采用式（13）。

语义检索和 SAM3 传播都是训练阶段的监督扩展机制。测试时，最终 U-Net 只接收一张图像并直接输出分割掩码；它不检索锚点、不构造路径，也不需要人工提示。这样既可以在训练阶段利用较复杂的跨图像推理扩大监督，又能保持部署阶段的轻量和稳定。

\bibliographystyle{plainnat}
\bibliography{../related_work_papers/references}

\end{document}
